\section{The first-match partition}\label{sec:partition}

Each predicate below is tested only after every earlier row has failed.
Here \(s=\min(u,t)\) in the chart selected by \eqref{eq:old:2_3}.
\[
\begin{array}{c|l}
\text{row}&\text{predicate}\\ \hline
O_R&R\le Y^{5/3-\epsilon_0}\\
O_M&M\le Y^{2/3-\epsilon_0}\\
O_{\operatorname{rad}M}&\operatorname{rad}(M)\le Y^{2/3-k_0}\\
Q0&K_{\rm res}>Y^{\Delta_{\rm exp}}\\
Q1&L<Y^{1/3-l_0}\\
Q2u&u\le t,\ \min(A,D)<Y^{1/3-j_0}\\
Q2t&t<u,\ \min(B,C)<Y^{1/3-j_0}\\
3u&u\le t,\ u\le Y^{a_0}\\
3t&t<u,\ t\le Y^{a_0}\\
4E&Y^{a_0}<s<Y^{b_0},\ \min(v,w)\le Y^{p_0}\\
4H&Y^{a_0}<s<Y^{b_0},\ v,w>Y^{p_0},\ \Pi\le Y^{p_0}\\
4P&Y^{a_0}<s<Y^{b_0},\ v,w>Y^{p_0},\ \Pi>Y^{p_0}\\
5W&\text{the remaining source.}
\end{array}\AutoTag{eq:old:3_1}
\]

\begin{theorem}[Exact partition]\label{thm:partition}
The thirteen first-match sets in \eqref{eq:old:3_1} are pairwise disjoint and have union
\(\Omega_Y\). On \(5W\),
\[
s\ge Y^{b_0},\qquad K_{\rm res}\le Y^{\Delta_{\rm exp}},\qquad
L\ge Y^{1/3-l_0},\qquad J,K\ge Y^{1/3-j_0}.
\AutoTag{eq:old:3_2}
\]
\end{theorem}

\begin{proof}
The quantities in \eqref{eq:old:3_1} are deterministic functions of the canonical
source. Each of the first twelve rows is intersected with the complements of
the preceding predicates. So these rows are disjoint. Their complement is
exactly the last row, which proves exhaustiveness. Failure of \(Q0\) gives
\(K_{\rm res}\le Y^{\Delta_{\rm exp}}\). Failure of \(Q1\) gives
\(L\ge Y^{1/3-l_0}\). On the complement of both chart-specific \(Q2\)
tests, the selected corners satisfy \(J,K\ge Y^{1/3-j_0}\).
Failure of the \(3u\) and \(3t\) tests gives \(s>Y^{a_0}\).
If \(s<Y^{b_0}\), trichotomy gives exactly one of the \(4E,4H,4P\)
predicates. Therefore their complement has \(s\ge Y^{b_0}\).
\end{proof}

The following boundary conventions apply after all earlier rows have
failed. The first three rows use non-strict upper cuts. Equality at
\(K_{\rm res}=Y^{\Delta_{\rm exp}}\),
\(L=Y^{1/3-l_0}\), and \(\min(J,K)=Y^{1/3-j_0}\) survives the
corresponding strict test. The chart equality \(u=t\) selects the
\(u\)-chart. The full first-match order determines its row. The
\(3u,3t\) tests include \(s=Y^{a_0}\), while \(5W\) contains
\(s=Y^{b_0}\). In the open middle interval, \(4E\) contains the boundary
\(\min(v,w)=Y^{p_0}\). Also \(4H\) contains \(\Pi=Y^{p_0}\) when
\(v,w>Y^{p_0}\).

\begin{remark}[Counting populations]
The estimates below count the following populations:
\begin{enumerate}
\item For \(O_R\), fix \(n\) first. Then apply the
      support-fibre lemma in \cite{Paper1} to the partners \(m\).
\item \(O_M,O_{\operatorname{rad}M},Q0\) use the complete fixed-\(M\)
      quadratic fibre, including all ordered factorizations
      \(M=(ab)(cd)\). \(Q1\) instead fixes \((P,L)\), with
      \(P=ab\) and \(Q_{\rm prod}=cd=P+L\).
\item \(Q2u,Q2t\) are disjoint chart populations. Their two analytic
      contributions apiece are additive bounds for each population.
\item \(3u,3t,4E,4H,4P\) retain the corresponding earlier-failure
      predicates and the canonical inverse filters. The \(4P\) estimate is
      not applied to the \(5W\) complement.
\item \(5W\) is exactly the remaining first-match set in \eqref{eq:old:3_1}.
\end{enumerate}
\end{remark}

Fix the ordered residual factors \((a,b,c,d)\), both shifts \((u,t)\),
the selected chart and its conic coordinates. Then the inverse identities
\eqref{eq:old:2_2} reconstruct the remaining two corners. This is the
full canonical label used for inverse reconstruction. The outer-row
divisor bound includes ordered residual factorizations. The joint-row sums
include their coefficient factorizations. The later estimates retain the
stated divisor or factor multiplicities. A larger set used to estimate one
row may intersect the set used for another row. The first-match partition
makes the original populations disjoint. Each row is bounded separately.
